\section{Stage 1A Eight Controls Full Results}
\label{app:stage1a}

Stage~1A evaluated whether the endpoint ranking of LeWM could be explained by
random geometry alone, by comparing the trained encoder (C0) against seven
controls on the same Track~A 100-pair artifact (8{,}000 records).  All
controls use Euclidean distance as the cost; only the representation
changes.  Ten random seeds were used for stochastic controls (C1--C5);
single-seed or deterministic controls are noted.

\subsection{Control Definitions}
\label{app:control_defs}

\begin{itemize}
\item \textbf{C0}: Trained LeWM encoder, 192-dimensional latent, Euclidean
  distance.  This is the baseline under audit.
\item \textbf{C1}: Same LeWM latents after a random orthogonal rotation
  ($192\times 192$); preserves all pairwise distances exactly.
\item \textbf{C2}: Gaussian random projection of LeWM latents to dimension
  $m \in \{1,2,4,8,16,32,64,128,192\}$, with entries
  $\mathcal{N}(0,1)/\!\sqrt{m}$.
\item \textbf{C3}: Coordinate subset of LeWM latents (first $m$ dimensions),
  $m \in \{8,16,32,64\}$.
\item \textbf{C4}: Gaussian null---i.i.d.\ $\mathcal{N}(0,1)$ vectors of
  dimension 192, independent of observations.
\item \textbf{C5}: Independent row shuffle of LeWM latents---each latent
  dimension is permuted independently, destroying inter-dimension structure.
\item \textbf{C6}: Randomly initialized LeWM encoder (same ViT-tiny architecture,
  untrained weights), seed~0.
\item \textbf{C7}: DINOv2 ViT-B/14 features (86.6M parameters,
  768-dimensional), CLS token and mean-pooled patch tokens.
\end{itemize}

\subsection{Aggregate Results}
\label{app:stage1a_agg}

Table~\ref{tab:stage1a-summary} reports the aggregate Spearman correlation,
pairwise accuracy, per-pair mean~$\rho$, and false-elite rate for each
control.  The $\pm$ values are standard deviations over 10 random seeds.

\begin{table}[h]
\centering
\caption{Stage~1A aggregate endpoint metrics on the Track~A 100-pair
artifact.  Spearman and pairwise accuracy are computed over all 8{,}000
records.  Per-pair $\rho$ averages 100 within-pair Spearman values.
False-elite rate is the fraction of the top-30 candidates (by model cost)
that are not in the true top-30 (by $C_\text{real\_state}$).
$\pm$ denotes std over 10 seeds for stochastic controls.}
\label{tab:stage1a-summary}
\resizebox{\textwidth}{!}{%
\begin{tabular}{llcccc}
\toprule
Control & Description & Spearman & Pairwise acc.\ & Per-pair $\rho$ & False elite \\
\midrule
C0 & LeWM 192-d & 0.506 & 0.632 & 0.333 & 0.700 \\
C1 & Orthogonal 192-d & 0.506$\pm$0.000 & 0.632$\pm$0.000 & 0.333 & 0.700 \\
\midrule
C2 ($m\!=\!1$) & Gaussian proj.\ & 0.196$\pm$0.082 & 0.554$\pm$0.014 & 0.145 & 0.790 \\
C2 ($m\!=\!8$) & Gaussian proj.\ & 0.379$\pm$0.035 & 0.603$\pm$0.013 & 0.266 & 0.723 \\
C2 ($m\!=\!32$) & Gaussian proj.\ & 0.444$\pm$0.033 & 0.621$\pm$0.011 & 0.309 & 0.706 \\
C2 ($m\!=\!64$) & Gaussian proj.\ & 0.473$\pm$0.026 & 0.626$\pm$0.007 & 0.320 & 0.704 \\
C2 ($m\!=\!192$) & Gaussian proj.\ & 0.495$\pm$0.019 & 0.629$\pm$0.007 & 0.326 & 0.701 \\
\midrule
C3 ($m\!=\!64$) & Coord.\ subset & 0.487$\pm$0.033 & 0.629$\pm$0.011 & 0.327 & 0.702 \\
\midrule
C4 & Gaussian null & $-$0.006$\pm$0.012 & 0.498$\pm$0.003 & $-$0.005 & 0.838 \\
C5 & Shuffled latent & 0.004$\pm$0.006 & 0.500$\pm$0.003 & $-$0.000 & 0.837 \\
\midrule
C6 & Randomly initialized ViT & $-$0.200 & 0.475 & $-$0.066 & 0.904 \\
\midrule
C7 (CLS) & DINOv2 CLS & 0.238 & 0.582 & 0.218 & 0.747 \\
C7 (mean) & DINOv2 mean & 0.286 & 0.595 & 0.253 & 0.750 \\
\bottomrule
\end{tabular}%
}
\end{table}

Key observations:

\paragraph{C1 confirms distance preservation.}
The orthogonal transform reproduces C0 exactly (Spearman 0.506, zero
variance), verifying that the evaluation pipeline is rotation-invariant.

\paragraph{C4 and C5 are uninformative.}
The Gaussian null (C4: $-$0.006) and shuffled latent (C5: 0.004) produce
near-zero Spearman, confirming that the endpoint ranking signal requires
\emph{structured} latent geometry, not arbitrary high-dimensional distances.

\paragraph{C6 is worse than random.}
The randomly initialized ViT-tiny reaches Spearman $-$0.200, \emph{below} both null
controls.  This anti-correlation is analyzed in detail in
Appendix~\ref{app:c6_audit} and discussed in Section~5.1.

\paragraph{C7 underperforms LeWM.}
DINOv2 mean-pool (0.286) and CLS (0.238) both fall below trained LeWM
(0.506), confirming that generic visual features do not match task-trained
endpoint geometry on PushT.

\subsection{C2 Gaussian Projection Ladder}
\label{app:c2_ladder}

Table~\ref{tab:c2-ladder} reports the full C2 dimension ladder from $m=1$
to $m=192$.  Endpoint Spearman rises steeply from $m=1$ to $m=32$, then
plateaus.  At $m=64$, Gaussian projections recover 93.5\% of the full
192-dimensional Spearman (0.473 vs.\ 0.506).  The C3 coordinate subset at
$m=64$ (0.487) is comparable, indicating that the ranking signal is
distributed rather than concentrated in specific learned dimensions.

\begin{table}[h]
\centering
\caption{C2 Gaussian projection and C3 coordinate subset endpoint Spearman
across projection dimensions.  $\pm$ denotes std over 10 seeds.}
\label{tab:c2-ladder}
\begin{tabular}{lccc}
\toprule
Dimension $m$ & C2 Spearman & C3 Spearman & C2 / C0 (\%) \\
\midrule
1 & 0.196$\pm$0.082 & --- & 38.7 \\
2 & 0.235$\pm$0.049 & --- & 46.4 \\
4 & 0.323$\pm$0.047 & --- & 63.8 \\
8 & 0.379$\pm$0.035 & 0.367$\pm$0.029 & 74.9 \\
16 & 0.407$\pm$0.037 & 0.423$\pm$0.032 & 80.4 \\
32 & 0.444$\pm$0.033 & 0.455$\pm$0.027 & 87.7 \\
64 & 0.473$\pm$0.026 & 0.487$\pm$0.033 & 93.5 \\
128 & 0.477$\pm$0.019 & --- & 94.3 \\
192 & 0.495$\pm$0.019 & --- & 97.8 \\
\bottomrule
\end{tabular}
\end{table}

This ladder is an \emph{endpoint} result.  Section~4 and
Appendix~\ref{app:stage1b} show that the same dimension sweep produces a
very different pattern when measured inside the final candidate pool of CEM:
$\Rpool$ stays near zero across all dimensions even as $\Rendpoint$ rises.

\subsection{Per-Anchor Breakdowns}
\label{app:anchor_breakdown}

The Track~A pairs include three diagnostic anchors defined in
Appendix~\ref{app:track_a}: the invisible quadrant (16 all-fail,
strong-$\rho$ pairs), the sign-reversal cluster (21 negative-$\rho$ pairs),
and the latent-favorable subset (12 pairs from D0xR1 and D1xR0).
Table~\ref{tab:anchor-breakdown} reports per-anchor Spearman for C0 and
C2 at $m=64$.

\begin{table}[h]
\centering
\caption{Per-anchor endpoint Spearman for C0 (trained LeWM) and C2 Gaussian
projection at $m=64$.  C2 values are mean$\pm$std over 10 seeds.}
\label{tab:anchor-breakdown}
\begin{tabular}{lcc}
\toprule
Anchor subset & C0 Spearman & C2 ($m\!=\!64$) Spearman \\
\midrule
Invisible quadrant (16 pairs) & 0.462 & 0.323$\pm$0.090 \\
Sign-reversal cluster (21 pairs) & $-$0.034 & 0.012$\pm$0.110 \\
Latent-favorable (12 pairs) & 0.519 & 0.479$\pm$0.045 \\
\midrule
All 100 pairs & 0.506 & 0.473$\pm$0.026 \\
\bottomrule
\end{tabular}
\end{table}

The invisible quadrant retains positive Spearman under both C0 (0.462) and
projected C2 (0.323), confirming that endpoint geometry looks healthy for
these pairs even though CEM cannot find successful actions.  The
sign-reversal cluster shows near-zero or slightly negative Spearman under
both representations, consistent with the encoder producing systematically
inverted rankings for these high-rotation pairs.

The latent-favorable subset preserves 92.3\% of C0 Spearman under C2
projection (0.479 vs.\ 0.519), indicating that the latent-favorable signal
is robust to dimensionality reduction.

\subsection{Implication for the Audit}

Stage~1A ruled out the ``Strong~B'' hypothesis: random
architecture-matched geometry does not explain the endpoint ranking of LeWM.
The trained encoder produces task-relevant structure (C0: 0.506) from a
worse-than-random starting point (C6: $-$0.200), and null controls
(C4, C5) confirm that arbitrary distances carry no signal.  Random
projections of the \emph{learned} space preserve much of this signal, but
that is a property of the learned geometry, not evidence that random
geometry suffices.  This conclusion motivated the Stage~1B planning
experiments (Appendix~\ref{app:stage1b}), which tested whether the preserved
endpoint signal transfers to CEM planning behavior.
